% Cote 119: offprints, reports, courses, CV and letters (typed and printed part): archive cote 119, site page 272 (PDF page 273).
% typescript. Draft compared: none.
% AI-transcribed by Claude Opus 5.5 (claude-opus-5-5, adaptive thinking, effort high), Message Batches API, batch r119_full (msgbatch_01DwZNm5Y7YsKBWyAN4sFs8w), 2026-11-13; not yet checked by eye.
\scanpage{119/272}
Il y a deux voies d'accès au doctorat d'Etat en Mathématiques.
A côté de la voie actuelle, qui reçoit désormais la mention ``recherche'',
est créée une voie d'enseignement destinée à faire une place importante aux critères
pédagogiques ; ce doctorat d'enseignement n'est pas un refuge pour ceux qui sont
inaptes à la recherche mais atteste une vocation authentique.

Les deux voies d'accès au doctorat (mention recherche et mention ensei-%
gnement) comportent un tronc commun : la thèse de spécialité, travail scientifique
écrit faisant l'objet d'un exposé oral du candidat, qui doit présenter un caractère
d'originalité soit par des qualités d'exposition et de mise au point d'une théorie
mathématique connue $^{(\hand{1})}$, soit par la contribution de résultats nouveaux, un article
publié dans un périodique, une monographie sont des exemples de thèse de spécialité .

Le candidat désireux de soutenir une thèse de spécialité en Mathématiques
à la Faculté des Sciences d'Orsay doit en faire la demande au département de Mathé-%
matiques, qui désigne un jury formé de trois membres, au moins, dont au moins deux
mathématiciens du département de Mathématiques d'Orsay. Le doctorat de spécialité
en Mathématiques est décerné sur avis unanime des membres du jury. La décision du
jury est sans appel.

Pour pouvoir être titulaire d'une thèse de spécialité, il n'y a aucune
condition d'âge, de possession de diplômes scolaires ou universitaires, de nationalité.

La possession de la thèse de spécialité est une condition nécessaire pour
pouvoir exercer la fonction de maître-assistant ; elle est suffisante sous réserve
de qualification pédagogique appropriée.

Pour pouvoir être candidat au titre de docteur ès-Sciences mathématiques
(mention recherche), i l faut réunir les conditions suivantes :

a) être titulaire d'une thèse de spécialité en Mathématiques,

b) avoir fait un travail de recherche original en Mathématiques présentant
un intérêt indiscutable,

c) justifier d'une culture mathématique solide.

Le candidat à un doctorat ès-Sciences mathématique (mention enseignement)
doit être titulaire d'une thèse de spécialité et doit prouver ses qualités pédagogi-%
ques en enseignant, comme M.A. pendant trois ans au moins dans trois enseignements
différents, dont un an au moins en 1er cycle et un an au moins en 2e cycle.

(\hand{1}) Ceci devrait permettre de combler de regrettables lacunes dans la littérature
mathématique qui sert de base à la recherche.

\mspage{272}
